% Phần 3 -- Yêu cầu 2: khái quát trước kiểu né lọc mới. Mọi số lấy qua \so{} (so_lieu.tex).
\section{Yêu cầu 2: khái quát trước kiểu né lọc mới}
\label{sec:yc2}

Nhóm đánh giá khả năng khái quát theo ba trục trên tập kiểm thử (Mục~\ref{sec:p1-test}): kiểu biến dạng chưa thấy
(C1, C2), kịch bản chưa thấy (\so{te.chuathay} đoạn C0) và giọng vùng miền. Kiểm tra rò rỉ giữa hai tập ghi nhận \so{rr.trunggan} cặp gần trùng (tương đồng 3-gram $\ge \so{rr.nguong}$) và \so{rr.trainkichban} đoạn
huấn luyện thuộc các kịch bản giữ ngoài.

C1 tác động lên chính tả, âm tiết nên có thể làm mất thông tin (``đánh'' và ``dành'' bỏ dấu đều thành ``danh''); C2
giữ nguyên âm tiết nhưng đổi lớp ký tự Unicode. Nhờ nhánh CNN đọc theo byte, v2 đạt Binary-F1 \so{C2.v2.bin} trên C2,
nhỉnh hơn C1 (\so{C1.v2.bin}); Macro-F1 trên hai tập gần như nhau (\so{C1.v2.macro} và \so{C2.v2.macro}).

\subsection{Theo từng kiểu biến dạng}

\begin{figure}[H]
\centering
% Sinh tự động bằng `python -m src.v2.eval.xuat_so_lieu_bao_cao`. Không sửa tay.
\begin{tikzpicture}
\draw[black!12] (0.00,0.3) -- (0.00,-7.80);
\node[font=\scriptsize,anchor=north] at (0.00,-7.80) {40};
\draw[black!12] (1.64,0.3) -- (1.64,-7.80);
\node[font=\scriptsize,anchor=north] at (1.64,-7.80) {50};
\draw[black!12] (3.27,0.3) -- (3.27,-7.80);
\node[font=\scriptsize,anchor=north] at (3.27,-7.80) {60};
\draw[black!12] (4.91,0.3) -- (4.91,-7.80);
\node[font=\scriptsize,anchor=north] at (4.91,-7.80) {70};
\draw[black!12] (6.55,0.3) -- (6.55,-7.80);
\node[font=\scriptsize,anchor=north] at (6.55,-7.80) {80};
\draw[black!12] (8.18,0.3) -- (8.18,-7.80);
\node[font=\scriptsize,anchor=north] at (8.18,-7.80) {90};
\node[font=\small,anchor=north] at (4.50,-8.25) {Macro-F1 5 nhãn};
\node[anchor=west,font=\small\bfseries] at (-5.2,0.00) {C1: hỏng chính con chữ};
\node[anchor=east,font=\footnotesize] at (-0.15,-0.55) {Dính từ};
\draw[black!30, line width=1.2pt] (2.84,-0.55) -- (7.71,-0.55);
\fill[maullm] (2.84,-0.55) circle (3.2pt);
\fill[mauvhai] (7.71,-0.55) circle (3.2pt);
\node[anchor=west,font=\scriptsize] at (7.83,-0.55) {87,1};
\node[anchor=east,font=\scriptsize,text=black!70] at (2.72,-0.55) {57,4};
\node[anchor=east,font=\footnotesize] at (-0.15,-1.05) {Kéo dài ký tự};
\draw[black!30, line width=1.2pt] (2.48,-1.05) -- (6.71,-1.05);
\fill[maullm] (2.48,-1.05) circle (3.2pt);
\fill[mauvhai] (6.71,-1.05) circle (3.2pt);
\node[anchor=west,font=\scriptsize] at (6.83,-1.05) {81,0};
\node[anchor=east,font=\scriptsize,text=black!70] at (2.36,-1.05) {55,1};
\node[anchor=east,font=\footnotesize] at (-0.15,-1.55) {Bỏ dấu};
\draw[black!30, line width=1.2pt] (2.96,-1.55) -- (6.56,-1.55);
\fill[maullm] (2.96,-1.55) circle (3.2pt);
\fill[mauvhai] (6.56,-1.55) circle (3.2pt);
\node[anchor=west,font=\scriptsize] at (6.68,-1.55) {80,1};
\node[anchor=east,font=\scriptsize,text=black!70] at (2.84,-1.55) {58,1};
\node[anchor=east,font=\footnotesize] at (-0.15,-2.05) {Số thay chữ};
\draw[black!30, line width=1.2pt] (4.50,-2.05) -- (6.47,-2.05);
\fill[maullm] (4.50,-2.05) circle (3.2pt);
\fill[mauvhai] (6.47,-2.05) circle (3.2pt);
\node[anchor=west,font=\scriptsize] at (6.59,-2.05) {79,5};
\node[anchor=east,font=\scriptsize,text=black!70] at (4.38,-2.05) {67,5};
\node[anchor=east,font=\footnotesize] at (-0.15,-2.55) {Chèn ký tự};
\draw[black!30, line width=1.2pt] (1.70,-2.55) -- (6.01,-2.55);
\fill[maullm] (1.70,-2.55) circle (3.2pt);
\fill[mauvhai] (6.01,-2.55) circle (3.2pt);
\node[anchor=west,font=\scriptsize] at (6.13,-2.55) {76,7};
\node[anchor=east,font=\scriptsize,text=black!70] at (1.58,-2.55) {50,4};
\node[anchor=east,font=\footnotesize] at (-0.15,-3.05) {Gõ sai};
\draw[black!30, line width=1.2pt] (2.19,-3.05) -- (5.36,-3.05);
\fill[maullm] (2.19,-3.05) circle (3.2pt);
\fill[mauvhai] (5.36,-3.05) circle (3.2pt);
\node[anchor=west,font=\scriptsize] at (5.48,-3.05) {72,8};
\node[anchor=east,font=\scriptsize,text=black!70] at (2.07,-3.05) {53,4};
\node[anchor=east,font=\footnotesize] at (-0.15,-3.55) {Teencode};
\draw[black!30, line width=1.2pt] (2.62,-3.55) -- (5.06,-3.55);
\fill[maullm] (2.62,-3.55) circle (3.2pt);
\fill[mauvhai] (5.06,-3.55) circle (3.2pt);
\node[anchor=west,font=\scriptsize] at (5.18,-3.55) {70,9};
\node[anchor=east,font=\scriptsize,text=black!70] at (2.50,-3.55) {56,0};
\node[anchor=west,font=\small\bfseries] at (-5.2,-4.25) {C2: đổi lớp ký tự (chưa thấy khi train)};
\node[anchor=east,font=\footnotesize] at (-0.15,-4.80) {Chữ toàn khổ};
\draw[black!30, line width=1.2pt] (5.52,-4.80) -- (7.02,-4.80);
\fill[maullm] (5.52,-4.80) circle (3.2pt);
\fill[mauvhai] (7.02,-4.80) circle (3.2pt);
\node[anchor=west,font=\scriptsize] at (7.14,-4.80) {82,9};
\node[anchor=east,font=\scriptsize,text=black!70] at (5.40,-4.80) {73,7};
\node[anchor=east,font=\footnotesize] at (-0.15,-5.30) {Kéo dài kép};
\draw[black!30, line width=1.2pt] (4.42,-5.30) -- (6.94,-5.30);
\fill[maullm] (4.42,-5.30) circle (3.2pt);
\fill[mauvhai] (6.94,-5.30) circle (3.2pt);
\node[anchor=west,font=\scriptsize] at (7.06,-5.30) {82,4};
\node[anchor=east,font=\scriptsize,text=black!70] at (4.30,-5.30) {67,0};
\node[anchor=east,font=\footnotesize] at (-0.15,-5.80) {Dấu phân cách lạ};
\draw[black!30, line width=1.2pt] (2.59,-5.80) -- (6.24,-5.80);
\fill[maullm] (2.59,-5.80) circle (3.2pt);
\fill[mauvhai] (6.24,-5.80) circle (3.2pt);
\node[anchor=west,font=\scriptsize] at (6.36,-5.80) {78,2};
\node[anchor=east,font=\scriptsize,text=black!70] at (2.47,-5.80) {55,8};
\node[anchor=east,font=\footnotesize] at (-0.15,-6.30) {Chữ Cyrillic giống Latin};
\draw[black!30, line width=1.2pt] (5.30,-6.30) -- (6.23,-6.30);
\fill[maullm] (5.30,-6.30) circle (3.2pt);
\fill[mauvhai] (6.23,-6.30) circle (3.2pt);
\node[anchor=west,font=\scriptsize] at (6.35,-6.30) {78,1};
\node[anchor=east,font=\scriptsize,text=black!70] at (5.18,-6.30) {72,4};
\node[anchor=east,font=\footnotesize] at (-0.15,-6.80) {Ký tự vô hình};
\draw[black!30, line width=1.2pt] (2.92,-6.80) -- (5.30,-6.80);
\fill[maullm] (2.92,-6.80) circle (3.2pt);
\fill[mauvhai] (5.30,-6.80) circle (3.2pt);
\node[anchor=west,font=\scriptsize] at (5.42,-6.80) {72,4};
\node[anchor=east,font=\scriptsize,text=black!70] at (2.80,-6.80) {57,9};
\node[anchor=east,font=\footnotesize] at (-0.15,-7.30) {Chèn dấu cách};
\draw[black!30, line width=1.2pt] (0.83,-7.30) -- (4.43,-7.30);
\fill[maullm] (0.83,-7.30) circle (3.2pt);
\fill[mauvhai] (4.43,-7.30) circle (3.2pt);
\node[anchor=west,font=\scriptsize] at (4.55,-7.30) {67,1};
\node[anchor=east,font=\scriptsize,text=black!70] at (0.71,-7.30) {45,1};
\fill[mauvhai] (1.0,0.75) circle (3.2pt); \node[anchor=west,font=\small] at (1.15,0.75) {v2};\fill[maullm] (2.6,0.75) circle (3.2pt); \node[anchor=west,font=\small] at (2.75,0.75) {gpt-oss-120b (API)};
\end{tikzpicture}

\caption{Macro-F1 5 nhãn theo 13 kiểu biến dạng, sắp giảm dần theo điểm của SafeViChi v2: gpt-oss-120b (API) (cam) và
SafeViChi v2 (xanh).}
\label{fig:p3-kieu}
\end{figure}

\begin{table}[H]
\centering
\footnotesize
\setlength{\tabcolsep}{3pt}
\caption{Kết quả theo từng kiểu biến dạng trên tập kiểm thử (Macro-F1, Bin-F1, BC F1 theo ý trên thang 0--100; FPR tính
theo \%; $n$: số đoạn; v1 là SafeViChi v1). Cột \textbf{Train C0}: kiểu có xuất hiện tự nhiên trong tập huấn luyện gốc
hay không.}
\label{tab:p3-kieu}
\begin{tabularx}{\linewidth}{@{}Y c c c c c c c c c@{}}
\toprule
 & & & \textbf{v1} & \multicolumn{2}{c}{\textbf{gpt-oss-120b (API)}} & \multicolumn{4}{c}{\textbf{SafeViChi v2}} \\
\cmidrule(lr){4-4}\cmidrule(lr){5-6}\cmidrule(l){7-10}
\textbf{Kiểu biến dạng} & \textbf{Train C0} & \textbf{$n$} & \textbf{Bin-F1} & \textbf{Macro} & \textbf{Bin-F1} & \textbf{Macro} & \textbf{Bin-F1} & \textbf{FPR} & \textbf{BC} \\
\midrule
\multicolumn{10}{@{}l}{\textit{C1: biến dạng bề mặt}} \\
Bỏ dấu & có & \so{kieu.C1.bodau.n} & \so{kieu.C1.bodau.v1.bin} & \so{kieu.C1.bodau.llm.macro} & \so{kieu.C1.bodau.llm.bin} & \textbf{\so{kieu.C1.bodau.v2.macro}} & \textbf{\so{kieu.C1.bodau.v2.bin}} & \so{kieu.C1.bodau.v2.fpr} & \so{kieu.C1.bodau.v2.bcf} \\
Teencode & có & \so{kieu.C1.teencode.n} & \so{kieu.C1.teencode.v1.bin} & \so{kieu.C1.teencode.llm.macro} & \so{kieu.C1.teencode.llm.bin} & \textbf{\so{kieu.C1.teencode.v2.macro}} & \textbf{\so{kieu.C1.teencode.v2.bin}} & \so{kieu.C1.teencode.v2.fpr} & \so{kieu.C1.teencode.v2.bcf} \\
Gõ sai & có & \so{kieu.C1.gosai.n} & \so{kieu.C1.gosai.v1.bin} & \so{kieu.C1.gosai.llm.macro} & \so{kieu.C1.gosai.llm.bin} & \textbf{\so{kieu.C1.gosai.v2.macro}} & \textbf{\so{kieu.C1.gosai.v2.bin}} & \so{kieu.C1.gosai.v2.fpr} & \so{kieu.C1.gosai.v2.bcf} \\
Kéo dài ký tự & có & \so{kieu.C1.keodai.n} & \so{kieu.C1.keodai.v1.bin} & \so{kieu.C1.keodai.llm.macro} & \so{kieu.C1.keodai.llm.bin} & \textbf{\so{kieu.C1.keodai.v2.macro}} & \textbf{\so{kieu.C1.keodai.v2.bin}} & \so{kieu.C1.keodai.v2.fpr} & \so{kieu.C1.keodai.v2.bcf} \\
Chèn ký tự & không & \so{kieu.C1.chenkytu.n} & \so{kieu.C1.chenkytu.v1.bin} & \so{kieu.C1.chenkytu.llm.macro} & \so{kieu.C1.chenkytu.llm.bin} & \textbf{\so{kieu.C1.chenkytu.v2.macro}} & \textbf{\so{kieu.C1.chenkytu.v2.bin}} & \so{kieu.C1.chenkytu.v2.fpr} & \so{kieu.C1.chenkytu.v2.bcf} \\
Số thay chữ & không & \so{kieu.C1.sothaychu.n} & \so{kieu.C1.sothaychu.v1.bin} & \so{kieu.C1.sothaychu.llm.macro} & \so{kieu.C1.sothaychu.llm.bin} & \textbf{\so{kieu.C1.sothaychu.v2.macro}} & \textbf{\so{kieu.C1.sothaychu.v2.bin}} & \so{kieu.C1.sothaychu.v2.fpr} & \so{kieu.C1.sothaychu.v2.bcf} \\
Dính từ & không & \so{kieu.C1.dinhtu.n} & \so{kieu.C1.dinhtu.v1.bin} & \so{kieu.C1.dinhtu.llm.macro} & \so{kieu.C1.dinhtu.llm.bin} & \textbf{\so{kieu.C1.dinhtu.v2.macro}} & \textbf{\so{kieu.C1.dinhtu.v2.bin}} & \so{kieu.C1.dinhtu.v2.fpr} & \so{kieu.C1.dinhtu.v2.bcf} \\
\midrule
\multicolumn{10}{@{}l}{\textit{C2: đổi lớp ký tự Unicode (điểm mù, hoàn toàn ngoài huấn luyện)}} \\
Chữ toàn khổ {\scriptsize(fullwidth)} & không & \so{kieu.C2.fullwidth.n} & \so{kieu.C2.fullwidth.v1.bin} & \so{kieu.C2.fullwidth.llm.macro} & \so{kieu.C2.fullwidth.llm.bin} & \textbf{\so{kieu.C2.fullwidth.v2.macro}} & \textbf{\so{kieu.C2.fullwidth.v2.bin}} & \so{kieu.C2.fullwidth.v2.fpr} & \so{kieu.C2.fullwidth.v2.bcf} \\
Ký tự vô hình {\scriptsize(zero-width)} & không & \so{kieu.C2.zerowidth.n} & \so{kieu.C2.zerowidth.v1.bin} & \so{kieu.C2.zerowidth.llm.macro} & \so{kieu.C2.zerowidth.llm.bin} & \textbf{\so{kieu.C2.zerowidth.v2.macro}} & \textbf{\so{kieu.C2.zerowidth.v2.bin}} & \so{kieu.C2.zerowidth.v2.fpr} & \so{kieu.C2.zerowidth.v2.bcf} \\
Chèn dấu cách {\scriptsize(space injection)} & không & \so{kieu.C2.spaceinjection.n} & \so{kieu.C2.spaceinjection.v1.bin} & \so{kieu.C2.spaceinjection.llm.macro} & \so{kieu.C2.spaceinjection.llm.bin} & \textbf{\so{kieu.C2.spaceinjection.v2.macro}} & \textbf{\so{kieu.C2.spaceinjection.v2.bin}} & \so{kieu.C2.spaceinjection.v2.fpr} & \so{kieu.C2.spaceinjection.v2.bcf} \\
Dấu phân cách lạ {\scriptsize(exotic separator)} & không & \so{kieu.C2.exoticseparator.n} & \so{kieu.C2.exoticseparator.v1.bin} & \so{kieu.C2.exoticseparator.llm.macro} & \so{kieu.C2.exoticseparator.llm.bin} & \textbf{\so{kieu.C2.exoticseparator.v2.macro}} & \textbf{\so{kieu.C2.exoticseparator.v2.bin}} & \so{kieu.C2.exoticseparator.v2.fpr} & \so{kieu.C2.exoticseparator.v2.bcf} \\
Chữ Cyrillic giống Latin {\scriptsize(homoglyph)} & không & \so{kieu.C2.cyrillichomoglyph.n} & \so{kieu.C2.cyrillichomoglyph.v1.bin} & \so{kieu.C2.cyrillichomoglyph.llm.macro} & \so{kieu.C2.cyrillichomoglyph.llm.bin} & \textbf{\so{kieu.C2.cyrillichomoglyph.v2.macro}} & \textbf{\so{kieu.C2.cyrillichomoglyph.v2.bin}} & \so{kieu.C2.cyrillichomoglyph.v2.fpr} & \so{kieu.C2.cyrillichomoglyph.v2.bcf} \\
Kéo dài kép {\scriptsize(double elongation)} & không & \so{kieu.C2.doubleelongation.n} & \so{kieu.C2.doubleelongation.v1.bin} & \so{kieu.C2.doubleelongation.llm.macro} & \so{kieu.C2.doubleelongation.llm.bin} & \textbf{\so{kieu.C2.doubleelongation.v2.macro}} & \textbf{\so{kieu.C2.doubleelongation.v2.bin}} & \so{kieu.C2.doubleelongation.v2.fpr} & \so{kieu.C2.doubleelongation.v2.bcf} \\
\bottomrule
\end{tabularx}
\end{table}

\textbf{Nhận xét:}
\begin{itemize}
  \item \textbf{SafeViChi v2 đứng đầu ở cả 13 kiểu biến dạng} về Macro-F1 và Binary-F1, kể cả kiểu chưa có
        khi huấn luyện như \emph{dính từ} (\so{kieu.C1.dinhtu.v2.macro} so với \so{kieu.C1.dinhtu.llm.macro} của LLM).
        Nhóm C1 giữ ngoài tập huấn luyện gốc còn đạt Macro-F1 cao hơn nhóm đã thấy (\so{nhom.C1.chuathay.v2.macro} so
        với \so{nhom.C1.dathay.v2.macro}).
  \item \textbf{Bằng chứng trên biến dạng điểm mù:} với \emph{chữ toàn khổ} và \emph{ký tự vô hình}, BC F1 của
        gpt-oss-120b (API) chỉ còn \so{kieu.C2.fullwidth.llm.bcf} và \so{kieu.C2.zerowidth.llm.bcf}, so với
        \so{kieu.C2.fullwidth.v2.bcf} và \so{kieu.C2.zerowidth.v2.bcf} của v2.
  \item \textbf{FPR tăng ở các kiểu phá vỡ ranh giới từ.} Phần lớn các kiểu giữ FPR thấp (\emph{bỏ dấu}
        \so{kieu.C1.bodau.v2.fpr}, \emph{chữ toàn khổ} \so{kieu.C2.fullwidth.v2.fpr}); FPR tăng chủ yếu ở \emph{chèn dấu cách} (\so{kieu.C2.spaceinjection.v2.fpr}), \emph{gõ
        sai} (\so{kieu.C1.gosai.v2.fpr}) và \emph{ký tự vô hình} (\so{kieu.C2.zerowidth.v2.fpr}). Do mỗi kiểu chỉ có
        khoảng \so{kieu.C2.spaceinjection.nat} đoạn \nhan{AN_TOAN}, một đoạn sai đã làm FPR đổi hơn 3 điểm phần trăm.
\end{itemize}

\subsection{Khái quát trên kịch bản chưa thấy}

\begin{figure}[H]
\centering
% Sinh tự động bằng `python -m src.v2.eval.xuat_so_lieu_bao_cao`. Không sửa tay.
\begin{tikzpicture}
\begin{axis}[ybar=2pt, bar width=13pt, width=\linewidth, height=6cm, ymin=0, ymax=100,
  symbolic x coords={C0,C1,C2}, xtick=data, enlarge x limits=0.22, ylabel={Macro-F1 5 nhãn},
  axis x line*=bottom, axis y line*=left, ymajorgrids, grid style={draw=black!12},
  tick label style={font=\small}, label style={font=\small},
  nodes near coords={\pgfmathprintnumber[fixed,fixed zerofill,precision=1,use comma]{\pgfplotspointmeta}},
  every node near coord/.append style={font=\tiny, color=black!75}, legend cell align=left,
  legend style={at={(0.5,1.03)}, anchor=south, legend columns=2, draw=none, font=\footnotesize,
    /tikz/every even column/.append style={column sep=0.5cm}}]
\addplot[fill=mauvhai, draw=white, area legend] coordinates {(C0,93.0) (C1,86.7) (C2,86.0)};
\addlegendentry{v2, kịch bản đã có trong train}
\addplot[fill=mauvhai!35, draw=mauvhai, area legend] coordinates {(C0,62.1) (C1,54.4) (C2,53.2)};
\addlegendentry{v2, kịch bản chưa thấy}
\addplot[fill=maullm, draw=white, area legend] coordinates {(C0,84.8) (C1,63.3) (C2,67.6)};
\addlegendentry{gpt-oss-120b, kịch bản đã có trong train v2}
\addplot[fill=maullm!35, draw=maullm, area legend] coordinates {(C0,53.0) (C1,39.0) (C2,45.5)};
\addlegendentry{gpt-oss-120b, kịch bản chưa thấy}
\end{axis}
\end{tikzpicture}

\caption{Macro-F1 5 nhãn trên kịch bản đã có trong tập huấn luyện (màu đậm) và kịch bản chưa thấy (màu nhạt).}
\label{fig:p3-kichban}
\end{figure}

\begin{table}[H]
\centering
\small
\setlength{\tabcolsep}{4.5pt}
\caption{Kết quả trên kịch bản trong phân phối huấn luyện và kịch bản chưa thấy (Macro-F1, Bin-F1, BC F1 theo ý
trên thang 0--100; FPR tính theo \%; $n$: số đoạn).}
\label{tab:p3-kichban}
\begin{tabular}{@{}l l c c c c c c@{}}
\toprule
 & & & \multicolumn{4}{c}{\textbf{SafeViChi v2}} & \textbf{gpt-oss-120b (API)} \\
\cmidrule(lr){4-7}\cmidrule(l){8-8}
\textbf{Tập} & \textbf{Phân nhóm} & \textbf{$n$} & \textbf{Macro-F1} & \textbf{Bin-F1} & \textbf{FPR} & \textbf{BC F1} & \textbf{Macro-F1} \\
\midrule
\multirow{2}{*}{C0} & Trong phân phối & \so{phan.C0.trongphanphoi.n} & \textbf{\so{phan.C0.trongphanphoi.v2.macro}} & \so{phan.C0.trongphanphoi.v2.bin} & \so{phan.C0.trongphanphoi.v2.fpr} & \so{phan.C0.trongphanphoi.v2.bcf} & \so{phan.C0.trongphanphoi.llm.macro} \\
 & Chưa thấy & \so{phan.C0.chuathay.n} & \textbf{\so{phan.C0.chuathay.v2.macro}} & \so{phan.C0.chuathay.v2.bin} & \so{phan.C0.chuathay.v2.fpr} & \so{phan.C0.chuathay.v2.bcf} & \so{phan.C0.chuathay.llm.macro} \\
\midrule
\multirow{2}{*}{C1} & Trong phân phối & \so{phan.C1.trongphanphoi.n} & \textbf{\so{phan.C1.trongphanphoi.v2.macro}} & \so{phan.C1.trongphanphoi.v2.bin} & \so{phan.C1.trongphanphoi.v2.fpr} & \so{phan.C1.trongphanphoi.v2.bcf} & \so{phan.C1.trongphanphoi.llm.macro} \\
 & Chưa thấy & \so{phan.C1.chuathay.n} & \textbf{\so{phan.C1.chuathay.v2.macro}} & \so{phan.C1.chuathay.v2.bin} & \so{phan.C1.chuathay.v2.fpr} & \so{phan.C1.chuathay.v2.bcf} & \so{phan.C1.chuathay.llm.macro} \\
\midrule
\multirow{2}{*}{C2} & Trong phân phối & \so{phan.C2.trongphanphoi.n} & \textbf{\so{phan.C2.trongphanphoi.v2.macro}} & \so{phan.C2.trongphanphoi.v2.bin} & \so{phan.C2.trongphanphoi.v2.fpr} & \so{phan.C2.trongphanphoi.v2.bcf} & \so{phan.C2.trongphanphoi.llm.macro} \\
 & Chưa thấy & \so{phan.C2.chuathay.n} & \textbf{\so{phan.C2.chuathay.v2.macro}} & \so{phan.C2.chuathay.v2.bin} & \so{phan.C2.chuathay.v2.fpr} & \so{phan.C2.chuathay.v2.bcf} & \so{phan.C2.chuathay.llm.macro} \\
\bottomrule
\end{tabular}
\end{table}

Trên C0, Macro-F1 của v2 giảm từ \so{phan.C0.trongphanphoi.v2.macro} (kịch bản đã thấy) xuống
\so{phan.C0.chuathay.v2.macro} (chưa thấy) nhưng vẫn cao hơn gpt-oss-120b (API) ở nhóm chưa thấy trên cả ba tập
(Bảng~\ref{tab:p3-kichban}). gpt-oss-120b (API) cũng giảm mạnh ở nhóm này (\so{phan.C0.trongphanphoi.llm.macro} $\to$
\so{phan.C0.chuathay.llm.macro}), cho thấy tình huống ẩn ý như đe dọa trả thù (DD6) hay kiểm soát người yêu (QR6)
vốn khó phân biệt giữa \nhan{DE_DOA}, \nhan{QUAY_ROI} và \nhan{AN_TOAN}. Nhóm chưa thấy ở C0 chỉ có \so{te.chuathay.xp}
đoạn \nhan{XUC_PHAM} nên Macro-F1 dao động mạnh theo từng mẫu.

\subsection{Theo giọng vùng miền}

\begin{table}[H]
\centering
\footnotesize
\setlength{\tabcolsep}{3.5pt}
\caption{Kết quả theo giọng vùng miền trên C0 (Macro-F1, Bin-F1 trên thang 0--100; FPR tính theo \%), kèm FPR của SafeViChi v2
trên C1 và C2. Chênh lệch là hiệu giữa vùng cao nhất và vùng thấp nhất.}
\label{tab:p3-vung}
\begin{tabular}{@{}l c c c c c c c c c@{}}
\toprule
 & & \multicolumn{3}{c}{\textbf{SafeViChi v2 (C0)}} & \multicolumn{3}{c}{\textbf{gpt-oss-120b (API, C0)}} & \multicolumn{2}{c}{\textbf{FPR v2}} \\
\cmidrule(lr){3-5}\cmidrule(lr){6-8}\cmidrule(l){9-10}
\textbf{Giọng} & \textbf{$n$} & \textbf{Macro} & \textbf{Bin-F1} & \textbf{FPR} & \textbf{Macro} & \textbf{Bin-F1} & \textbf{FPR} & \textbf{C1} & \textbf{C2} \\
\midrule
Bắc & \so{C0.bac.v2.n} & \textbf{\so{C0.bac.v2.macro}} & \textbf{\so{C0.bac.v2.bin}} & \so{C0.bac.v2.fpr} & \so{C0.bac.llm.macro} & \so{C0.bac.llm.bin} & \so{C0.bac.llm.fpr} & \so{C1.bac.v2.fpr} & \so{C2.bac.v2.fpr} \\
Trung & \so{C0.trung.v2.n} & \textbf{\so{C0.trung.v2.macro}} & \textbf{\so{C0.trung.v2.bin}} & \so{C0.trung.v2.fpr} & \so{C0.trung.llm.macro} & \so{C0.trung.llm.bin} & \so{C0.trung.llm.fpr} & \so{C1.trung.v2.fpr} & \so{C2.trung.v2.fpr} \\
Nam & \so{C0.nam.v2.n} & \textbf{\so{C0.nam.v2.macro}} & \textbf{\so{C0.nam.v2.bin}} & \so{C0.nam.v2.fpr} & \so{C0.nam.llm.macro} & \so{C0.nam.llm.bin} & \so{C0.nam.llm.fpr} & \so{C1.nam.v2.fpr} & \so{C2.nam.v2.fpr} \\
\midrule
Chênh lệch & & \textbf{\so{C0.v2.gap.macro}} & \textbf{\so{C0.v2.gap.bin}} & \so{C0.v2.gap.fpr} & \so{C0.llm.gap.macro} & \so{C0.llm.gap.bin} & \so{C0.llm.gap.fpr} & \so{C1.v2.gap.fpr} & \so{C2.v2.gap.fpr} \\
\bottomrule
\end{tabular}
\end{table}

SafeViChi v2 có Macro-F1 và Binary-F1 cao hơn gpt-oss-120b (API) ở cả ba giọng (Bảng~\ref{tab:p3-vung}). Chênh lệch
Macro-F1 giữa các vùng của v2 là \so{C0.v2.gap.macro} điểm, nhỏ hơn \so{C0.llm.gap.macro} của gpt-oss-120b (API).
Ngược lại, FPR của v2 chênh lớn giữa các vùng (Trung \so{C0.trung.v2.fpr}, Bắc \so{C0.bac.v2.fpr}, Nam
\so{C0.nam.v2.fpr}; chênh \so{C0.v2.gap.fpr}) và lặp lại trên C1, C2, trong khi FPR của gpt-oss-120b (API) gần như
bằng nhau (chênh \so{C0.llm.gap.fpr}). Tuy vậy, mỗi vùng chỉ có
\so{C0.nam.v2.at.n}--\so{C0.bac.v2.at.n} đoạn \nhan{AN_TOAN} nên mỗi đoạn sai làm FPR đổi khoảng 3 điểm phần trăm;
cần thêm mẫu an toàn theo vùng trước khi kết luận v2 thiên lệch theo giọng.

\subsection{Các lát cắt khác}

\begin{table}[H]
\centering
\small
\caption{Kết quả trên một số lát cắt ngữ cảnh, độ lộ và kênh giao tiếp ở tập kiểm thử C0 (thang 0--100; $n$: số đoạn).}
\label{tab:p3-latcat}
\begin{tabular}{@{}l c c c c@{}}
\toprule
 & & \multicolumn{2}{c}{\textbf{SafeViChi v2}} & \textbf{gpt-oss-120b (API)} \\
\cmidrule(lr){3-4}\cmidrule(l){5-5}
\textbf{Lát cắt} & \textbf{$n$} & \textbf{Macro-F1} & \textbf{Bin-F1} & \textbf{Macro-F1} \\
\midrule
Hội thoại 1 dòng & \so{luot.C0.1dong.n} & \textbf{\so{luot.C0.1dong.v2.macro}} & \so{luot.C0.1dong.v2.bin} & \so{luot.C0.1dong.llm.macro} \\
Hội thoại 2--8 dòng & \so{luot.C0.nhieudong.n} & \textbf{\so{luot.C0.nhieudong.v2.macro}} & \so{luot.C0.nhieudong.v2.bin} & \so{luot.C0.nhieudong.llm.macro} \\
Đoạn nguy hại lộ liễu & \so{dolo.C0.lolieu.n} & \textbf{\so{dolo.C0.lolieu.v2.macro}} & \so{dolo.C0.lolieu.v2.bin} & \so{dolo.C0.lolieu.llm.macro} \\
Đoạn nguy hại ẩn ý & \so{dolo.C0.any.n} & \textbf{\so{dolo.C0.any.v2.macro}} & \so{dolo.C0.any.v2.bin} & \so{dolo.C0.any.llm.macro} \\
Kênh Messenger/Zalo & \so{kenh.C0.messengerzalo.n} & \textbf{\so{kenh.C0.messengerzalo.v2.macro}} & \so{kenh.C0.messengerzalo.v2.bin} & \so{kenh.C0.messengerzalo.llm.macro} \\
Kênh nhóm lớp & \so{kenh.C0.nhomlop.n} & \textbf{\so{kenh.C0.nhomlop.v2.macro}} & \so{kenh.C0.nhomlop.v2.bin} & \so{kenh.C0.nhomlop.llm.macro} \\
Kênh SMS & \so{kenh.C0.sms.n} & \so{kenh.C0.sms.v2.macro} & \so{kenh.C0.sms.v2.bin} & \textbf{\so{kenh.C0.sms.llm.macro}} \\
\bottomrule
\end{tabular}
\end{table}

Theo Bảng~\ref{tab:p3-latcat}, khoảng cách giữa v2 và gpt-oss-120b (API) lớn nhất ở hội thoại 2--8 dòng và đoạn nguy hại
ẩn ý. Kênh SMS (\so{kenh.C0.sms.n} đoạn, nhiều tin một dòng thiếu ngữ cảnh như thông báo OTP, giao hàng hợp lệ) là lát
cắt duy nhất gpt-oss-120b (API) cao hơn.

\subsection{Phân tích một số trường hợp thất bại}

\begin{table}[H]
\centering
\footnotesize
\caption{Bốn đoạn nguy hại thuộc kịch bản chưa thấy bị SafeViChi v2 dự đoán là \nhan{AN_TOAN} trên C0 (mỗi dòng trong ô là một
lượt nói).}
\label{tab:p3-vidu}
\begin{tabularx}{\linewidth}{@{}>{\raggedright\arraybackslash}p{0.41\linewidth} >{\raggedright\arraybackslash}p{0.20\linewidth} Y@{}}
\toprule
\textbf{Đoạn tin nhắn} & \textbf{Kết quả} & \textbf{Nguyên nhân khó} \\
\midrule
\so{vd.ghiam.text} &
Nhãn chuẩn: \so{vd.ghiam.gold}\newline v2: \so{vd.ghiam.pred}\newline LLM (API): \so{vd.ghiam.llm} &
Kịch bản \so{vd.ghiam.kichban} (đe dọa trả thù vì bị báo thầy cô). Đe dọa ẩn trong câu ``nhớ đi đường đông người nhé'', không có từ thô tục; gpt-oss-120b (API) cũng bỏ sót. \\
\midrule
\so{vd.vay.text} &
Nhãn chuẩn: \so{vd.vay.gold}\newline v2: \so{vd.vay.pred}\newline LLM (API): \so{vd.vay.llm} &
Kịch bản \so{vd.vay.kichban} (vay trực tuyến chỉ cần căn cước) giữ ngoài tập huấn luyện. Đòi giấy tờ thay vì chuyển tiền ngay nên khác mẫu lừa đảo quen thuộc. \\
\midrule
\so{vd.nham.text} &
Nhãn chuẩn: \so{vd.nham.gold}\newline v2: \so{vd.nham.pred}\newline LLM (API): \so{vd.nham.llm} &
Chiêu ``chuyển nhầm tiền'' hoàn toàn mới, lời lẽ lịch sự như nhờ giúp đỡ thông thường. \\
\midrule
\so{vd.kiemsoat.text} &
Nhãn chuẩn: \so{vd.kiemsoat.gold}\newline v2: \so{vd.kiemsoat.pred}\newline LLM (API): \so{vd.kiemsoat.llm} &
Kịch bản \so{vd.kiemsoat.kichban} (kiểm soát người yêu) giữ ngoài tập huấn luyện. Không có từ xúc phạm hay đe dọa thân thể; nguy hại nằm ở việc ép kiểm soát tài khoản và đòi gửi ảnh chụp màn hình. \\
\bottomrule
\end{tabularx}
\end{table}

Điểm chung của bốn trường hợp: mô hình học tốt mẫu đã thấy nhưng chưa nắm khái niệm tổng quát về ``ép buộc''
hay ``dụ chuyển tiền'' khi cách diễn đạt hoàn toàn mới.
